\documentclass[10pt,a4paper]{article}
\usepackage[brazil]{babel}
\usepackage[utf8]{inputenc}
\renewcommand{\baselinestretch}{1.0}
\usepackage{mhchem}
\usepackage{graphicx}
\usepackage{svg}
\usepackage{multicol}
\usepackage{amssymb}
\usepackage{amsmath}
\usepackage{enumitem}
\usepackage{geometry}
\geometry{a4paper, margin=2cm}
\setlist{nosep, topsep=2pt, partopsep=0pt, parsep=0pt, itemsep=3pt}

\newcommand{\oC}{$^\circ$C}
\newcommand{\e}[1]{$\times 10^{#1}$}
\sffamily

\newcommand{\questaoheader}[1]{%
  \begin{center}
  {\small\sffamily QG101 --- Bloco 9: Forças Intermoleculares e Estados da Matéria
  \hfill Questão #1}
  \end{center}
  \vspace{-0.3cm}
  \noindent\rule{\textwidth}{0.6pt}
  \vspace{0.1cm}
}

\newcommand{\resolucaosep}{%
  \vspace{0.4cm}
  \noindent\rule{\textwidth}{0.4pt}
  \vspace{0.1cm}

  \noindent{\bf Resolução:}
  \vspace{0.2cm}
}

\begin{document}
\pagestyle{empty}

% ================================================================
% Questão 1
% ================================================================
\questaoheader{1}

\noindent\textbf{Questão 1.}
Coloque as substâncias \ce{CH4}, \ce{H2O}, \ce{NH3} e \ce{Ne} em ordem
\textbf{crescente} de ponto de ebulição e indique, para cada uma, o
principal tipo de força intermolecular envolvido.

\resolucaosep

\noindent Ordem crescente de ponto de ebulição:
$$\ce{Ne} < \ce{CH4} < \ce{NH3} < \ce{H2O}$$

\begin{itemize}
  \item \ce{Ne} (PE $= -246$\oC): átomo nobre, muito pequeno e de baixa
        polarizabilidade --- \textbf{dispersão de London} (forças muito fracas).
  \item \ce{CH4} (PE $= -161$\oC): molécula apolar e de massa molar pequena ---
        \textbf{dispersão de London} (mais forte que no Ne por ter mais
        elétrons e maior polarizabilidade).
  \item \ce{NH3} (PE $= -33$\oC): possui H ligado a N (muito eletronegativo) ---
        \textbf{ligações de hidrogênio}, consideravelmente mais fortes que as
        forças de London do \ce{CH4}.
  \item \ce{H2O} (PE $= 100$\oC): O é ainda mais eletronegativo que N e cada
        molécula pode formar até 4 ligações de hidrogênio (2 doando H, 2
        aceitando) --- \textbf{ligações de hidrogênio} (as mais fortes do
        conjunto).
\end{itemize}

\clearpage

% ================================================================
% Questão 2
% ================================================================
\questaoheader{2}

\noindent\textbf{Questão 2.}
A tabela abaixo apresenta os pontos de ebulição dos hidretos do grupo 16:

\begin{center}
\begin{tabular}{lcccc}
\hline
 & \ce{H2S} & \ce{H2Se} & \ce{H2Te} & \ce{H2O} \\
\hline
P.E. (\oC) & $-60$ & $-41$ & $-2$ & $100$ \\
\hline
\end{tabular}
\end{center}

\begin{enumerate}
  \item Explique por que o ponto de ebulição cresce de \ce{H2S} para
        \ce{H2Te}.
  \item Se a água seguisse a mesma tendência (apenas forças de London),
        seu ponto de ebulição seria de aproximadamente $-90$\oC. Explique
        por que o valor real é tão maior.
\end{enumerate}

\resolucaosep

\begin{enumerate}
  \item De \ce{H2S} para \ce{H2Te} as moléculas ficam progressivamente
        maiores: o número de elétrons cresce (16 → 34 → 52), aumentando a
        \textbf{polarizabilidade eletrônica}. Dipólos instantâneos mais
        intensos geram forças de dispersão de London mais fortes, exigindo
        mais energia para vaporizar e, portanto, elevando o ponto de
        ebulição.

  \item A água apresenta \textbf{ligações de hidrogênio} muito fortes:
        o oxigênio é altamente eletronegativo ($\chi = 3{,}44$) e de raio
        atômico pequeno, tornando o vínculo O--H muito polar. Cada molécula
        pode formar até 4 ligações H (2 como doadora, 2 como aceitadora),
        criando uma rede coesiva intensa. Para romper essas interações e
        vaporizar a água é necessária energia muito superior à prevista
        pelas forças de London sozinhas, elevando o PE real para 100\oC
        em vez de $\approx -90$\oC.
\end{enumerate}

\clearpage

% ================================================================
% Questão 3
% ================================================================
\questaoheader{3}

\noindent\textbf{Questão 3.}
Para cada uma das espécies abaixo, identifique o \textbf{principal} tipo
de força intermolecular presente no estado líquido/sólido:
\begin{enumerate}
  \item \ce{HF}
  \item \ce{CO2}
  \item \ce{CH4}
  \item \ce{NH3}
  \item \ce{I2}
\end{enumerate}

\resolucaosep

\begin{enumerate}
  \item \ce{HF}: \textbf{ligação de hidrogênio} (H ligado ao F, o elemento
        mais eletronegativo; forças H muito intensas).
  \item \ce{CO2}: \textbf{dispersão de London} (molécula linear e simétrica
        --- apolar; os dipolos C=O se cancelam).
  \item \ce{CH4}: \textbf{dispersão de London} (molécula apolar de baixo
        peso molecular).
  \item \ce{NH3}: \textbf{ligação de hidrogênio} (H ligado ao N
        eletronegativo; também há par solitário no N que aceita H de outra
        molécula).
  \item \ce{I2}: \textbf{dispersão de London} (molécula apolar, porém
        grande e com muitos elétrons, resultando em forças de London
        relativamente fortes --- \ce{I2} é sólido a 25\oC).
\end{enumerate}

\clearpage

% ================================================================
% Questão 4
% ================================================================
\questaoheader{4}

\noindent\textbf{Questão 4.}
Explique, em termos de estrutura e de ligações de hidrogênio, por que o
gelo flutua na água líquida. Por que essa propriedade é considerada
anômala em relação à maioria das substâncias?

\resolucaosep

No gelo, cada molécula de água forma \textbf{4 ligações de hidrogênio}
com vizinhas (2 doando H, 2 aceitando), organizando-se numa
\textbf{rede hexagonal aberta} (estrutura cristalina com amplos canais
vazios). Essa geometria fixa faz o gelo ocupar \emph{mais} volume por
molécula do que a água líquida, resultando em \textbf{menor densidade}
do sólido ($\approx 0{,}92$~g/cm$^3$) em relação ao líquido
($\approx 1{,}00$~g/cm$^3$).

Quando o gelo funde, parte das ligações H se rompe, permitindo que as
moléculas se aproximem e se empacotem de forma mais compacta --- daí a
densidade aumentar ao fundir.

A anomalia: na \textbf{maioria das substâncias} o sólido é mais denso
que o líquido (o empacotamento ordenado do cristal gera maior densidade),
portanto o sólido afunda no próprio líquido. A água inverte esse
comportamento por causa da estrutura aberta do gelo imposta pelas
ligações de hidrogênio direcionais.

\clearpage

% ================================================================
% Questão 5
% ================================================================
\questaoheader{5}

\noindent\textbf{Questão 5.}
O etanol (\ce{C2H5OH}) é completamente solúvel em água, mas o hexano
(\ce{C6H14}) praticamente não se dissolve. Explique essa diferença usando
o princípio ``semelhante dissolve semelhante'' e os tipos de força
intermolecular envolvidos em cada caso.

\resolucaosep

\begin{itemize}
  \item \textbf{Etanol em água}: o etanol possui o grupo hidroxila (--OH),
        que forma \textbf{ligações de hidrogênio} com as moléculas de água.
        As interações soluto--solvente (etanol--água) são do mesmo tipo das
        interações solvente--solvente (água--água), tornando a dissolução
        energeticamente favorável. Como ambos são polares e capazes de
        estabelecer ligações H, eles se misturam em qualquer proporção
        (\emph{semelhante dissolve semelhante}).

  \item \textbf{Hexano em água}: o hexano é completamente apolar e só
        apresenta \textbf{forças de dispersão de London}. Ao tentar se
        dissolver em água, seria necessário romper a forte rede de ligações
        de hidrogênio da água sem substituí-la por interações igualmente
        fortes com o hexano --- processo energeticamente desfavorável. As
        moléculas de água se ``excluem'' do hexano, que permanece segregado.
\end{itemize}

Em resumo: polar dissolve polar (etanol em água); apolar não dissolve em
polar (hexano em água).

\clearpage

% ================================================================
% Questão 6
% ================================================================
\questaoheader{6}

\noindent\textbf{Questão 6.}
O butano (\ce{C4H10}, $M = 58$~g/mol) ferve a $-1$\oC, enquanto a acetona
(\ce{C3H6O}, $M = 58$~g/mol) ferve a $56$\oC.
\begin{enumerate}
  \item As duas substâncias têm a mesma massa molar. O que explica a
        grande diferença nos pontos de ebulição?
  \item Quais forças intermoleculares atuam em cada uma?
\end{enumerate}

\resolucaosep

\begin{enumerate}
  \item A diferença se deve à \textbf{polaridade molecular}: o butano é
        apolar, enquanto a acetona possui uma ligação \ce{C=O} (grupo
        carbonila) com grande diferença de eletronegatividade, gerando um
        \textbf{dipolo permanente} intenso. Forças dipolo--dipolo
        adicionam-se às forças de London na acetona, exigindo
        $\approx 57$\oC a mais para vaporizar, mesmo com a mesma massa
        molar.

  \item
    \begin{itemize}
      \item Butano (\ce{C4H10}): apenas \textbf{dispersão de London}
            (molécula apolar).
      \item Acetona (\ce{C3H6O}): \textbf{dispersão de London} +
            \textbf{interações dipolo--dipolo} (molécula polar com dipolo
            permanente no grupo C=O).
    \end{itemize}
\end{enumerate}

\clearpage

% ================================================================
% Questão 7
% ================================================================
\questaoheader{7}

\noindent\textbf{Questão 7.}
À temperatura ambiente, \ce{F2} e \ce{Cl2} são gases, \ce{Br2} é líquido e
\ce{I2} é sólido, embora todos sejam moléculas apolares formadas pelo
mesmo tipo de ligação covalente simples.
\begin{enumerate}
  \item Qual força intermolecular é responsável pelas diferenças de
        estado físico observadas?
  \item Relacione essa tendência ao número de elétrons e à
        polarizabilidade de cada molécula.
\end{enumerate}

\resolucaosep

\begin{enumerate}
  \item A força responsável é a \textbf{dispersão de London} (forças de
        London ou forças de dispersão de van der Waals). Todas as moléculas
        são apolares, portanto não há dipolo permanente; as atrações
        surgem de \emph{dipolos instantâneos} induzidos por flutuações
        eletrônicas.

  \item À medida que se desce no grupo 17, o número de elétrons cresce
        drasticamente (\ce{F2}: 18 e$^-$; \ce{Cl2}: 34 e$^-$; \ce{Br2}:
        70 e$^-$; \ce{I2}: 106 e$^-$). Mais elétrons em orbitais mais
        difusos significa maior \textbf{polarizabilidade}: a nuvem
        eletrônica se deforma mais facilmente, gerando dipolos instantâneos
        maiores e interações de London mais intensas. Por isso:
        \begin{itemize}
          \item \ce{F2} e \ce{Cl2}: forças de London fracas → gases.
          \item \ce{Br2}: forças intermediárias → líquido.
          \item \ce{I2}: forças de London fortes → sólido.
        \end{itemize}
\end{enumerate}

\clearpage

% ================================================================
% Questão 8
% ================================================================
\questaoheader{8}

\noindent\textbf{Questão 8.}
A glicerina (com três grupos \ce{-OH} por molécula) é muito mais viscosa
que a água, que por sua vez é mais viscosa que a gasolina.
\begin{enumerate}
  \item Relacione a viscosidade de cada líquido com o tipo e a quantidade
        de forças intermoleculares presentes.
  \item Por que a água apresenta tensão superficial mais alta que a
        maioria dos solventes orgânicos comuns?
\end{enumerate}

\resolucaosep

\begin{enumerate}
  \item A viscosidade reflete a resistência das moléculas ao fluxo, que
        depende da intensidade das interações entre elas:
        \begin{itemize}
          \item \textbf{Glicerina} (mais viscosa): 3 grupos --OH por
                molécula formam uma rede extensa de \textbf{ligações de
                hidrogênio} ($\sim$6--9 ligações H por molécula, contando
                doação e aceitação), dificultando muito o movimento relativo
                das moléculas.
          \item \textbf{Água} (viscosidade intermediária): até 4 ligações H
                por molécula --- fortes, mas menos numerosas que na glicerina.
          \item \textbf{Gasolina} (menos viscosa): mistura de
                hidrocarbonetos apolares, apenas forças de London (fracas) →
                moléculas deslizam com facilidade.
        \end{itemize}

  \item A tensão superficial reflete o custo energético de criar nova
        superfície. Na água, as moléculas da superfície perderam parte de
        sua rede de \textbf{ligações de hidrogênio} (para baixo e para os
        lados, mas não para cima), resultando num estado de alta energia. O
        sistema minimiza essa energia contraindo a superfície, o que se
        manifesta como alta tensão superficial. Solventes orgânicos comuns
        possuem apenas forças de London ou dipolo--dipolo (mais fracas),
        logo a diferença energética entre moléculas de superfície e de
        interior é menor --- tensão superficial menor.
\end{enumerate}

\clearpage

% ================================================================
% Questão 9
% ================================================================
\questaoheader{9}

\noindent\textbf{Questão 9.}
Quando o \ce{NaCl} se dissolve em água, os íons \ce{Na+} e \ce{Cl-} ficam
rodeados por moléculas de água (hidratação).
\begin{enumerate}
  \item Que tipo de interação (distinta das três discutidas neste bloco)
        ocorre entre os íons e as moléculas de água?
  \item As moléculas de água se orientam de forma específica ao redor de
        cada íon. Indique qual extremidade da molécula de água
        (\ce{O} ou \ce{H}) fica voltada para o \ce{Na+} e qual fica
        voltada para o \ce{Cl-}, justificando.
\end{enumerate}

\resolucaosep

\begin{enumerate}
  \item O tipo de interação é a \textbf{força íon--dipolo}: uma carga
        iônica (total) interage com o dipolo permanente da molécula de água
        (polar). Essa força é geralmente mais intensa que as ligações de
        hidrogênio e é a principal responsável pela solvatação de eletrólitos
        em água.

  \item A molécula de água tem o oxigênio com carga parcial negativa
        ($\delta^-$) e os hidrogênios com carga parcial positiva ($\delta^+$):
        \begin{itemize}
          \item Ao redor do \ce{Na+} (carga $+$): a extremidade
                \textbf{oxigênio} ($\delta^-$) fica voltada para o cátion,
                pois cargas opostas se atraem.
          \item Ao redor do \ce{Cl-} (carga $-$): as extremidades
                \textbf{hidrogênio} ($\delta^+$) ficam voltadas para o
                ânion, pela mesma razão de atração eletrostática.
        \end{itemize}
\end{enumerate}

\clearpage

% ================================================================
% Questão 10
% ================================================================
\questaoheader{10}

\noindent\textbf{Questão 10.}
A figura abaixo mostra a curva de aquecimento de uma amostra de água pura
sob pressão constante, partindo do estado sólido.

\begin{center}
\includesvg[width=0.55\textwidth]{figuras/bloco09_fig2_curva_aquecimento}
\end{center}

\begin{enumerate}
  \item Por que a temperatura permanece constante durante os trechos de
        fusão e de ebulição, mesmo com fornecimento contínuo de calor?
  \item Para onde vai a energia fornecida durante esses trechos
        constantes, já que ela não aumenta a temperatura?
  \item O trecho de ebulição é mais longo (no eixo do calor) que o trecho
        de fusão. O que isso indica sobre as energias envolvidas em cada
        mudança de estado?
\end{enumerate}

\resolucaosep

\begin{enumerate}
  \item A temperatura mede a energia cinética média das moléculas. Durante
        uma transição de fase, a energia fornecida é usada para
        \textbf{romper forças intermoleculares} (e não para acelerar as
        moléculas), portanto a energia cinética --- e a temperatura ---
        permanecem constantes até que toda a amostra tenha mudado de fase.

  \item A energia vai para \textbf{vencer as atrações intermoleculares}:
        \begin{itemize}
          \item Na fusão: rompe-se parcialmente a rede de ligações de
                hidrogênio do gelo (estrutura cristalina) para que as
                moléculas possam se mover livremente no líquido.
          \item Na ebulição: rompem-se praticamente \emph{todas} as
                interações restantes da fase líquida, permitindo que as
                moléculas escapem para a fase gasosa de forma independente.
        \end{itemize}

  \item O trecho de ebulição mais longo indica que a
        \textbf{entalpia de vaporização} ($\Delta H_{\text{vap}}
        \approx 2260$~J/g) é muito maior que a
        \textbf{entalpia de fusão} ($\Delta H_{\text{fus}}
        \approx 334$~J/g). Isso faz sentido: ao fundir, apenas parte da
        rede de ligações H é rompida (o líquido ainda mantém muitas
        interações); ao vaporizar, é necessário vencer quase todas as
        forças intermoleculares restantes para separar completamente as
        moléculas.
\end{enumerate}

\end{document}
